\section{Frame laws and preparation of unit laws}\label{sec:frame-laws}

We now prepare the probability laws used to prove raw supersaturation.
Under the reference law \(\mu^2\), prescribed vector images have
exponentially small probabilities, even when their nominal directions
combine both endpoints of a unit. This will imply two useful properties
of any law satisfying \Cref{eq:raw-law-caps}. After discarding
exponentially small mass, tensors shared by its endpoints have sparse
profiles, and the law can be partitioned into
parts on which every unpinned exact-image constraint remains unlikely.
The normalized laws of these parts will be called leaf laws.

\subsection{Uniform frame laws}

We record the linear algebra behind all reference image laws.
For a subspace of a nominal coefficient space, its \emph{image} is its image under the relevant frame map.
Nominal independence is independence before applying these maps.

\begin{lemma}\label{lem:paired-frame-orbit}
Let \(A,A':U\to V\) and \(B,B':W\to V^*\) be injective maps between finite-dimensional spaces over \(\F\).
If
\[
 B\trans A=(B')\trans A',
\]
there is \(L\in\GL(V)\) with \(LA=A'\) and \(L^{-\mathsf T}B=B'\).
Thus based individually injective plus and minus frames with a prescribed Gram matrix form one orbit under the simultaneous primal and dual action.
\end{lemma}
\begin{proof}
Write \(\Phi=B\trans:V\to W^*\) and \(\Phi'=(B')\trans:V\to W^*\).
Both are surjective.
The based map \(f:A(U)\to A'(U)\) given by \(f(Au)=A'u\) satisfies \(\Phi'f=\Phi\).
It therefore maps \(A(U)\cap\ker\Phi\) isomorphically onto \(A'(U)\cap\ker\Phi'\).
Extend this restriction to an isomorphism \(L_0:\ker\Phi\to\ker\Phi'\).

Choose a complement \(U_0\) of \(A(U)\cap\ker\Phi\) in \(A(U)\).
Its images under \(\Phi\) are independent.
Extend a basis of \(\Phi(U_0)\) to a basis of \(W^*\), and choose lifts of the added basis vectors under \(\Phi\) and \(\Phi'\), respectively.
For the basis vectors in \(\Phi(U_0)\), choose the lifts in \(U_0\) and their images under \(f\).
Define \(L\) to equal \(L_0\) on \(\ker\Phi\) and to carry each chosen lift to its primed lift.
These are direct-sum decompositions of \(V\), so \(L\) is an isomorphism.
It extends \(f\) and satisfies \(\Phi'L=\Phi\), which is equivalent to the two required frame identities.
\end{proof}

\begin{lemma}[Gram normalization]\label{lem:gram-frame-normalization}
Let \(p,q\) be nonnegative integers with \(p+q\le N\).
For \(p\) independent uniform plus vectors and \(q\) independent uniform minus vectors in paired copies of \(\F^N\), and a prescribed \(p\)-by-\(q\) Gram matrix \(G\), the probability of that Gram matrix together with individual injectivity of both frames is at most \(2^{-pq}\) and at least
\begin{equation}\label{eq:bounded-gram-normalization}
 2^{-pq}(1-2^{p-N})(1-2^{p+q-N}).
\end{equation}
In particular, any consistent Gram specification on a bounded number of vectors, together with injectivity on each side, has probability bounded below by a positive constant for large \(N\).
\end{lemma}
\begin{proof}
The plus columns are linearly independent with probability at least \(1-2^{p-N}\): each nonzero coefficient dependence has probability \(2^{-N}\).
Given their linear independence, each minus column takes its required evaluations with probability \(2^{-p}\), independently, so the Gram costs \(2^{-pq}\).
Under this Gram conditioning, every nonzero combination of minus columns is uniform on an affine space of dimension at least \(N-p\).
Its probability of vanishing is at most \(2^{-(N-p)}\).
Union over fewer than \(2^q\) combinations gives failure probability at most \(2^{p+q-N}\).
This proves the lower bound.
For the upper bound with plus injectivity required, condition on the plus columns and use the exact \(2^{-pq}\) Gram probability whenever they are independent.
A partially specified consistent Gram contains some full specification; for bounded \(p,q\), the same positive lower bound for that full specification proves the last assertion.
\end{proof}

\begin{lemma}\label{lem:frame-images}
The following statements hold for all sufficiently large \(n\).
\begin{enumerate}[label=\textup{(\alph*)}]
\item
The plus frame \([P,X]\) is uniform among injective maps.
Conditional on it, the minus columns are independent uniforms on the affine spaces prescribed in \Cref{eq:raw-frame-law}, followed only by the conditioning that \([Q,Y]\) is injective.
The failure probability of this last condition, before conditioning, is at most \(2^{2(d+h)-N}\) per component.

Given \(P,Q\), the channels \(X,Y\) are independent uniform-column draws on \(\ker Q\trans\) and \(\ker P\trans\), respectively, conditioned on \([P,X]\) and \([Q,Y]\) being injective.
The two rank-failure probabilities are exponentially small in \(N\).
Each channel frame alone is uniform among injective maps \(\F^h\to\F^N\).
For any fixed \(q\) independent vectors in its paired ambient space, the probability that transpose evaluation on their span is not injective is at most
\begin{equation}\label{eq:channel-transpose-failure}
 2^{q-h+1}.
\end{equation}
The unconditional joint channel law is bounded above by a constant, depending on \(g,h\) but not on \(n\), times the completely independent uniform-column law.

\item
At independently sampled raw vertices and components, images of prescribed primal coefficient tuples that are independent separately at each vertex and sign are uniform subject to their individual injectivities and their prescribed same-vertex Gram matrices.
For tuples using channels as well, the same assertion holds conditionally on the full self-channel Gram matrices \(X_e\trans Y_e\); the relevant full Gram matrix is
\(\begin{psmallmatrix}E&0\\0&X_e\trans Y_e\end{psmallmatrix}\).

\item
Let a fixed tuple of nominal directions use arbitrary combinations of the two endpoints' primal and channel coefficients.
Compute its rank separately in each component and sign in the \emph{direct sum over both endpoints}, and let the total be \(t\).
For every preassigned \(\epsilon>0\), choosing \(M_0\) sufficiently large gives, under \(\mu^2\),
\begin{equation}\label{eq:reference-image-cap}
 \PP\{\text{all those images equal prescribed vectors}\}
 \le 2^{-(1-\epsilon)tN}\qquad(t\ge1).
\end{equation}
This is valid for all possible ranks \(t\), including \(t=O(n)\), not only for bounded tuples.
\end{enumerate}
\end{lemma}
\begin{proof}
The raw law is invariant under \(L\) on each \(V_e^+\) and \(L^{-\mathsf T}\) on \(V_e^-\), independently at vertices and components.
For (a), the action is transitive on injective plus frames, so their marginal is uniform.
Given a plus frame, each column of \(Q\) has specified evaluations against \(P,X\), while each column of \(Y\) has zero evaluations against \(P\).
There is no other constraint until full minus injectivity is imposed.
The remaining allowed set is therefore the product of those affine column spaces restricted to full rank.

Each column space contains the common translation space
\(\operatorname{ann}(\im[P,X])\), of dimension \(N-d-h\).
For a nonzero coefficient vector \(c\in\F^{d+h}\), the corresponding linear combination of minus columns is consequently uniform on an affine space containing that translation space.
Its chance of being zero is at most \(2^{-(N-d-h)}\).
Union over the fewer than \(2^{d+h}\) nonzero coefficient vectors gives the asserted bound \(2^{2(d+h)-N}\).

Alternatively, condition first on \(P,Q\).
The allowed \(X\)-columns lie in \(\ker Q\trans\), the allowed \(Y\)-columns in \(\ker P\trans\), and their two full-rank restrictions are separate.
This proves the stated conditional independence.
For example, before imposing injectivity of \([P,X]\), any dependence with a nonzero \(X\)-coefficient has probability at most \(2^{-(N-d)}\); union over its at most \(2^{d+h}\) coefficient choices gives a bound \(2^{2d+h-N}\).
The estimate for \(Y\) is identical.

The marginal of \(X\) alone is invariant under \(\GL(\F^N)\) and supported on injective maps; transitivity gives its uniformity, and likewise for \(Y\).
Before conditioning independent channel columns to be injective, their evaluations on a fixed independent \(q\)-tuple form a uniform \(h\)-by-\(q\) matrix.
For each nonzero coefficient vector in \(\F^q\), the chance that its evaluation is zero is \(2^{-h}\).
Union over those coefficient vectors gives at most \(2^{q-h}\), and channel injectivity has probability at least \(1/2\) for large \(N\).
This proves \Cref{eq:channel-transpose-failure}.
To compare their joint marginal with completely independent columns, condition on \(G=X\trans Y\).
The conditional marginal is uniform on the orbit of two individually injective \(h\)-frames having Gram matrix \(G\).
Under independent uniform columns that orbit has probability at least \(2^{-h^2-2}\) for large \(N\), by \Cref{lem:gram-frame-normalization}.
Since the raw probability of any \(G\) is at most one, the density ratio on every such orbit is at most \(2^{h^2+2}\).
Multiply over the fixed number of components to obtain the joint channel bound.
Its possible dependence on \(h\) is harmless: \(h\) is fixed before \(n\) grows.

For (b), the restricted Gram of primal tuples is fixed by \(E\).
By \Cref{lem:paired-frame-orbit}, all based realizations with that Gram and individual injectivity lie in one orbit.
Projection of an invariant uniform raw law is uniform on this orbit: an element carrying one partial realization to another bijects their full-frame completions.
There is at least one completion because these tuples are restrictions of a raw frame.
After fixing \(X_e\trans Y_e\), the same argument applies to arbitrary full-frame coefficient tuples.
Independence over raw draws and components proves (b).

For (c), we prove the image cap, keeping explicit the dependence on the rank.
Replace each tested tuple by a basis of its nominal span; inconsistent dependent prescriptions have probability zero.
For the plus directions at the two endpoints, start with completely independent uniform plus columns and then impose individual full-frame injectivity.
The image of a rank-\(t_+\) tuple in the joint nominal direct sum is exactly uniform on \(t_+\) independent ambient vectors before this conditioning.
Indeed, at each ambient coordinate the coefficient map is a surjective linear map onto \(\F^{t_+}\), and different ambient coordinates use independent input bits.
Full injectivity has probability \(1-2^{-\Omega(N)}\), uniformly over the given construction, so the plus prescription costs at most
\((1-2^{-\Omega(N)})^{-1}2^{-t_+N}\).

Now condition on all plus frames.
At a component let \(S\) be the span of both endpoints' plus columns; then \(\dim S\le2(d+h)\).
Before minus injectivity, every minus column has an affine distribution whose translation space contains \(S^\perp\).
Split off an independent uniform \(S^\perp\)-part of each column and condition on the remaining parts.
For a rank-\(t_-\) tuple of minus coefficient combinations, the \(S^\perp\)-parts of its images are independent uniform vectors in \(S^\perp\), because the coefficient map has rank \(t_-\).
Thus prescribed minus images have probability at most
\[
 2^{-(N-2(d+h))t_-}.
\]
The additional conditioning on the two minus frames being injective changes this by a factor \(1+2^{-\Omega(N)}\), by (a).
Perform this argument componentwise.
There is a fixed total conditioning factor bounded, for example, by \(4^{|\cE|}\), and the resulting bound is
\[
 4^{|\cE|}\,2^{-t_+N-(N-2(d+h))t_-}.
\]
Choose \(M_0\) so that \(2(d+h)/N<\epsilon/2\) for large \(n\), and then increase \(n\) so that the fixed prefactor is at most \(2^{\epsilon N/2}\).
Since \(t=t_++t_-\ge1\), this proves \Cref{eq:reference-image-cap} uniformly for every allowable \(t\).
In particular, mixing the two endpoints in a direction causes no loss of an entire \(N\)-bit image constraint.
\end{proof}

\subsection{Sparse profiles arising from true intersections}

For a matrix profile, its \emph{total component rank} means \(\sum_{e\in\cE}\rank x_e\).

\begin{lemma}\label{lem:raw-intersections}
Let \(\sigma\) satisfy the joint cap in \Cref{eq:raw-law-caps}.
Outside a set of \(\sigma\)-probability \(2^{-\Omega(N)}\),
\[
 \sum_{e\in\cE}
  \dim\bigl(\im P_{1,e}\cap\im P_{2,e}\bigr)\le2D.
\]
For any unit outside that exceptional set and any
\(U_1\lambda_1=U_2\lambda_2\), both \(\lambda_i\) have total component rank at most \(2D\).
Each has a unique representation \((w_{i,l})_l\) vanishing at a strict majority of tags, and
\[
 |\{l:w_{i,l}\ne0\}|\le4D/g,\qquad
 \chi_*(w_1)=\chi_*(w_2).
\]
\end{lemma}
\begin{proof}
If the summed intersection dimension exceeds \(2D\), choose \(2D\) independent common-image vectors, componentwise.
Write each as \(P_{1,e}a=P_{2,e}b\).
Their nominal relation directions \((a,b)\) are independent in the direct sum of the two primal coefficient spaces at that component.
Indeed a dependence between the pairs would, after applying \(P_{1,e}\), be the corresponding dependence between the selected common-image vectors.
Across components their ranks add.

There are at most \(2^{O(n)}\) ways to choose this bounded number of coefficient pairs and their component indices.
For each choice, \Cref{eq:reference-image-cap} bounds the probability of their zero images by \(2^{-(1-\epsilon)2DN}\) under \(\mu^2\).
After paying \(2^{DN}\) for \(\sigma\), and making the coefficient count a sufficiently small fraction of \(N\) by the choice of \(M_0\), the union bound is \(2^{-\Omega(N)}\).

Suppose \(U_1\lambda_1=U_2\lambda_2=\xi\).
The column space of \(\xi_e\) lies in both primal plus spans.
Since multiplication by the two injective frame maps preserves matrix rank,
\[
 \rank(\lambda_{i,e})=\rank(\xi_e)
 \le\dim(\im P_{1,e}\cap\im P_{2,e}).
\]
This proves the total-rank assertion.

Take any representation \((w_l)_l\) of one of the profiles.
A component \(w_l+w_{l'}\) is nonzero whenever \(w_l\ne w_{l'}\), so there are at most \(2D\) unequal unordered pairs.
If no value occurred at a strict majority of tags, write \(c_v\) for its multiplicities.
Then
\[
 \#\{\text{unequal pairs}\}
 =\frac12\left(g^2-\sum_vc_v^2\right)
 \ge\frac12\left(g^2-\frac g2\sum_vc_v\right)
 =\frac{g^2}{4}>2D,
\]
a contradiction.
Let the majority have size \(g-z\).
There are at least \(z(g-z)\ge zg/2\) unequal pairs, hence \(z\le4D/g\).
Its common value belongs to \(W_*\) by \Cref{lem:cut-kernel}, so subtract it from every coordinate.
The resulting representation vanishes at a strict majority.
It is unique: two such representations differ by a constant, and their two majority zero sets intersect.

Finally, equality of the component tensors gives equality of their ambient traces.
By cyclicity of trace and \Cref{eq:E-moments},
\[
 \Tr(P_{i,e}\lambda_{i,e}Q_{i,e}\trans)
 =\pair{P_{i,e}\trans Q_{i,e}}{\lambda_{i,e}}
 =\eta(\lambda_{i,e}).
\]
For \(e=\{l,l'\}\), it follows that
\[
 \eta(w_{1,l})+\eta(w_{2,l})
 =\eta(w_{1,l'})+\eta(w_{2,l'}).
\]
This value is independent of the tag.
The two sparse representations together have at most \(8D/g=32000<g\) nonzero coordinates, so some tag has both entries zero.
The common value is therefore zero.
Summing over tags proves equality of the two specified \(\chi_*\)-values.
\end{proof}

\subsection{Peeling into leaves with exact pins}

We partition almost all of a capped unit law into leaves by storing
unusually likely exact-image constraints in bounded-dimensional pin spaces.
The remaining directions then satisfy a uniform image-probability bound.
For a unit \(O=(1,2)\), introduce nominal coefficient spaces at each component
\[
 \cL_{O,e}
   =\bigoplus_{i\in O}(\cB_i\oplus H_i^+),
 \qquad
 \cR_{O,e}
   =\bigoplus_{i\in O}(\cB_i\oplus H_i^-),
 \qquad H_i^\pm=\F^h.
\]
Their frame-image maps add the endpoint images in the common ambient \(V_e^\pm\).
A \emph{pin space} is a specified subspace \(D_{O,e}^\pm\) of this nominal direct sum together with specified exact images of every vector in a chosen basis.
It may mix both endpoints, and it may mix primal and channel coordinates.
Once the basis images are known, every image on the pin space is fixed.

For a tested tuple of directions, its \emph{rank modulo the pins} is the sum, over components and signs, of the dimensions of the spans of its classes in the nominal quotients by \(D_{O,e}^\pm\).
The rank is nominal; it is not computed after applying the raw frames.

The budgets for this procedure are fixed before the selector, mixer,
and channel parameters in the construction. Their values are
\begin{equation}\label{eq:pin-budgets}
 \begin{gathered}
 k_{\max}=56(g-1),\qquad
 \zeta=\frac{1}{1000k_{\max}},\qquad
 K_1=\left\lceil\frac{4(D+1)}{\zeta}\right\rceil.
 \end{gathered}
\end{equation}
The phase and second-pinning budgets are
\begin{equation}\label{eq:law-budgets}
 \begin{gathered}
 r=2K_1+2,\qquad L_0=\lceil100(D+10)\rceil,\qquad
 d_0=4rL_0,\\
 u_0=10(K_1+d_0+1),\qquad
 K=\left\lceil\frac{10(D+K_1+u_0+10)}{\zeta}\right\rceil .
 \end{gathered}
\end{equation}
Here the integer \(r\) is a rank budget, distinct from the affine-gradient
map \(r(x)\). The complete order of choices is verified in
\Cref{sec:parameters}.

\begin{lemma}[Two peelings]\label{lem:peeling}
Take the constants in \Cref{eq:pin-budgets,eq:law-budgets}, and choose the reference accuracy in \Cref{eq:reference-image-cap} with \(\epsilon<\zeta/4\).
For every law \(\sigma\) satisfying \Cref{eq:raw-law-caps}, the following decompositions are available.
\begin{enumerate}[label=\textup{(\alph*)}]
\item
After discarding the exception in \Cref{lem:raw-intersections} and at most \(2^{-\zeta N}\) further mass, the law is partitioned into \emph{leaves}.
Each leaf has fixed pin spaces with total dimension at most \(K_1\), and its normalized law \(\sigma_\lambda\) satisfies
\begin{equation}\label{eq:leaf-minentropy}
 \sigma_\lambda\{\text{a specified exact-image constraint of rank }t
                 \text{ modulo the pins}\}
 \le 2^{-(1-\zeta)tN}
 \qquad(t\ge1).
\end{equation}
The normalized leaf laws satisfy a common bound \(2^{O(N)}\mu^2\), with the implied constant determined by the early parameters.

\item
Suppose a restriction of the mixed law has relative mass \(2^{-o(N)}\).
One may discard an exponentially small fraction of that restricted law so that every remaining normalized leaf restriction satisfies
\begin{equation}\label{eq:final-leaf-minentropy}
 \sigma_\lambda\{\text{a specified exact-image constraint of rank }t
                 \text{ modulo the pins}\}
 \le 2^{-(1-2\zeta)tN}
 \qquad(t\ge1),
\end{equation}
and still has density \(2^{O(N)}\) relative to \(\mu^2\).
For several successive restrictions, the same statement applies to their combined relative cost.

\item
In addition, suppose at most \(u_0\) further coefficient directions per unit must be pinned, their coefficient choices and all required scalar metadata having at most \(2^{O(n)}\) possibilities per old leaf.
After splitting by these choices and their exact images, discarding exponentially small mass, and peeling again, the total pin dimension is at most \(K\).
The fresh leaves satisfy \Cref{eq:leaf-minentropy}; after further restrictions as in \textup{(b)}, they satisfy \Cref{eq:final-leaf-minentropy}.
They retain a uniform \(2^{O(N)}\mu^2\) density bound.
\end{enumerate}
In all parts the mixed law uses each leaf's actual mass, normalized only after the indicated discards.
No equal reweighting of leaves is used.
After a restriction, \(\sigma_\lambda\) denotes the normalized surviving
law on that leaf.
\end{lemma}

The per-leaf density bounds in \Cref{lem:peeling} are joint bounds
relative to \(\mu^2\); they do not supply the original one-vertex
marginal caps. Restricting the \emph{mixed} law by mass \(2^{-o(N)}\),
on the other hand, gives marginal bounds \(M2^{o(N)}\mu\) and joint
bound \(2^{(D+o(1))N}\mu^2\). These are the two forms of control used
later: leaf estimates govern averaged key distributions, and mixed-law
estimates govern common tests. Their probability scales are summarized
in \Cref{sec:probability-scales}.

\begin{proof}
We give the finite procedure and its quantitative bounds.
Work first with the original, unnormalized restriction of \(\sigma\) obtained by removing the intersection exception.
In particular its density is still at most \(2^{DN}\) relative to \(\mu^2\).
Let \(R\) be its residual support.
As long as its residual mass exceeds \(2^{-\zeta N}\), begin with its normalized restriction and no pins.
If \Cref{eq:leaf-minentropy} fails, condition on a violating image constraint and add its nominal directions to the pins.
Given the existing pins, a consistent tuple constraint can be reduced to a basis of its span modulo the pins.
The other prescriptions are either automatic or inconsistent; an inconsistent event cannot violate a positive probability bound.
Thus every step adds positive rank.

If the cumulative added rank is \(u\), the unnormalized mass of the current part is greater than
\[
 2^{-\zeta N-(1-\zeta)uN}.
\]
The cumulative directions are independent after taking a basis, so the absolute density cap and \Cref{eq:reference-image-cap} give the upper bound
\[
 2^{DN-(1-\epsilon)uN}.
\]
Comparing the bounds yields
\begin{equation}\label{eq:peeling-rank-comparison}
 (\zeta-\epsilon)u<D+\zeta.
\end{equation}
This is strictly below the rank budget \(K_1\).
The same comparison applies to a whole violating tuple at once, so a large tuple cannot jump beyond that budget.
The rank therefore cannot increase indefinitely.
The process reaches a part with no violating constraint; remove that entire part as a leaf and restart on the residual.
The underlying raw unit space is finite, and every leaf is nonempty, so the procedure terminates with residual mass at most \(2^{-\zeta N}\).

A leaf of rank \(u\le K_1\) has mass greater than
\(2^{-\zeta N-(1-\zeta)uN}\).
Dividing the original density cap by this lower bound shows, for example,
\[
 \sigma_\lambda\le
 2^{(D+K_1+1)N}\mu^2.
\]
This proves (a), including uniformity of its density constant.

For (b), write the current normalized mixture as
\(\sum_\lambda w_\lambda\sigma_\lambda\).
Let \(q_\lambda\) be the relative mass retained by the combined restriction in leaf \(\lambda\), and let
\(q=\sum_\lambda w_\lambda q_\lambda=2^{-o(N)}\).
Discard leaves for which \(q_\lambda<2^{-\zeta N}\).
Their mass in the restricted law is at most
\[
 q^{-1}\sum_{\lambda:q_\lambda<2^{-\zeta N}}
       w_\lambda q_\lambda
 \le 2^{-\zeta N+o(N)}.
\]
On every other leaf, division by \(q_\lambda\) multiplies probabilities by at most \(2^{\zeta N}\).
For \(t\ge1\),
\[
 2^{\zeta N}2^{-(1-\zeta)tN}
 \le2^{-(1-2\zeta)tN}.
\]
This proves the claimed min-entropy and density bounds.
It also explains why the small mass bound is needed for the mixed law, rather than separately asserted for every old leaf.

For (c), first apply this trimming to any preceding small-cost restrictions.
Split each retained old leaf by the further coefficient choices, scalar metadata, and exact images.
There are at most \(2^{O(n)}2^{u_0N}\) parts.
Discard every part having relative mass less than
\(2^{-(u_0+1)N}\) in its old leaf.
The total discarded fraction is at most
\[
 2^{O(n)-N}=2^{-\Omega(N)}
\]
once \(M_0\) is large enough for the fixed metadata coefficient.
The normalized law \(Q\) on each retained part consequently has the absolute density bound
\begin{equation}\label{eq:second-peeling-start}
 Q\le2^{D'N}\mu^2,\qquad D'=D+K_1+u_0+5.
\end{equation}
The harmless slack covers the preceding trimming and normalizations.
There are initially at most \(K_1+u_0\) pinned dimensions on this part.

Repeat the first procedure, keeping those old and deliberately added pins fixed.
Only new directions are counted in its rank comparison.
To justify this point, let \(c_1,\ldots,c_u\) be cumulative new lifts independent modulo the stored pin space.
They are linearly independent in the full nominal space.
The refined event is contained in their prescribed image event, whether or not the old pin event is also imposed.
Using the \emph{unconditioned} reference measure in \Cref{eq:second-peeling-start},
\[
 Q(\text{refined event})
 \le 2^{D'N}
     \mu^2\{\operatorname{image}(c_j)\text{ has its prescribed value for all }j\}
 \le2^{D'N-(1-\epsilon)uN}.
\]
No reference measure conditioned on the rare old-pin event is used.
For each selected current part the tuple is fixed, so adaptive selection of the next violating tuple does not require an additional union bound.

The lower mass bound is again \(2^{-\zeta N-(1-\zeta)uN}\).
Thus \(u<4(D'+1)/\zeta\), a conservative consequence of the analogue of \Cref{eq:peeling-rank-comparison}.
Together with the initial \(K_1+u_0\) pins this is less than the chosen \(K\) in \Cref{eq:law-budgets}.
The same leaf-mass calculation gives a \(2^{O(N)}\mu^2\) bound with a common early-parameter constant.
The second procedure discards at most \(2^{-\zeta N}\) relative mass per part; summing with the actual part weights gives the same total bound.
It restores \Cref{eq:leaf-minentropy} relative to all old and new pins.
Part (b) can then be applied once more to yield \Cref{eq:final-leaf-minentropy} after subsequent restrictions.

Throughout these operations, disintegration over the finite parts is exact.
Multiplying each normalized part by its actual mass recovers the retained law.
This proves the last assertion as well.
\end{proof}

\subsection{Bilinear sign estimates}

\begin{lemma}[Walsh estimates]\label{lem:walsh}
Let all weights below have absolute value at most one.
\begin{enumerate}[label=\textup{(\alph*)}]
\item
If \(x,y\) are independent uniform vectors and \(B(x,y)\) is a bilinear form of rank \(d\), then
\[
 \left|\E f(x)g(y)(-1)^{B(x,y)}\right|\le2^{-d/2}.
\]
Replacing either uniform law by a law of density at most \(C_i\) multiplies this bound by at most \(C_1C_2\).

\item
If \(\alpha,\beta\) are subprobability measures on \(\F^d\), with point masses at most \(p_1,p_2\), respectively, then
\[
 \left|\sum_{x,y}\alpha(x)\beta(y)f(x)g(y)(-1)^{x\cdot y}\right|
 \le2^{d/2}(p_1p_2)^{1/2}.
\]

\item
For independent uniform global \(N\)-bit vector slots in two groups, a cross-pairing character with coefficient matrix \(C\) has bit rank \(N\rank C\).
In particular, if its cross coefficient pattern is nonzero, its correlation against two bounded weights, one on each group, is at most \(2^{-N/2}\).
Arbitrary phases depending separately on the groups may be included in these weights.
\end{enumerate}
\end{lemma}
\begin{proof}
Let \(H\) be the \(2^d\)-square matrix with entries \(H_{xy}=(-1)^{x\cdot y}\).
For \(x,x'\in\F^d\),
\[
 (HH\trans)_{xx'}=\sum_y(-1)^{(x+x')\cdot y}
 =\begin{cases}2^d,&x=x',\\0,&x\ne x'.\end{cases}
\]
The second value is zero by pairing \(y\) with \(y+e_j\) at any coordinate where \(x+x'\) is nonzero.
Thus the Euclidean operator norm of \(H\) is \(2^{d/2}\).

For (b), set \(a_x=\alpha(x)f(x)\) and \(b_y=\beta(y)g(y)\).
Then
\[
 |a\trans Hb|\le2^{d/2}\|a\|_2\|b\|_2,\qquad
 \|a\|_2^2\le p_1\sum_x\alpha(x)\le p_1,
\]
and similarly for \(b\), proving the result.

For (a), invertible changes of coordinates reduce the form to the dot product on its \(d\) active coordinates.
Average each bounded weight over its inactive coordinates.
Apply (b) to the uniform probabilities \(p_1=p_2=2^{-d}\).
A bounded-density change can be absorbed into the weights after dividing their bounds by \(C_1,C_2\).

Finally, the matrix of the full bit form in (c) is \(C\otimes I_N\).
Choose bases reducing \(C\) to \(\operatorname{diag}(I_{\rank C},0)\); the tensor matrix then has rank \(N\rank C\).
Apply (a), absorbing any within-group phases into the corresponding bounded weight.
\end{proof}
